\documentclass[11pt,letterpaper]{article}

% Page layout
\usepackage[margin=1in]{geometry}
\usepackage[T1]{fontenc}
\usepackage{lmodern}
\usepackage{microtype}

% Math and economics notation
\usepackage{amsmath,amssymb,amsthm,mathtools,bm}

% Figures and tables
\usepackage{graphicx}
\usepackage{booktabs,threeparttable}
\usepackage{xcolor}
\usepackage{float}
\usepackage{array,longtable}
\usepackage{caption,subcaption}

% Lists and numbers
\usepackage{enumitem}
\usepackage{siunitx}

% References and links
\usepackage{natbib}
\usepackage[hidelinks]{hyperref}
\usepackage[capitalise]{cleveref} % load after hyperref

% Optional theorem environments
\newtheorem{proposition}{Proposition}
\newtheorem{lemma}{Lemma}

\setlength{\parindent}{0pt}
\setlength{\parskip}{0.55em}

\begin{document}

\begin{center}
    {\LARGE\bfseries Lecture Notes}\\[6pt]
    {\large Your Name --- Course Name}\\[3pt]
\end{center}

\section*{1. Main Idea}
Explain the economic question and the key intuition.

\section*{2. Model and Reasoning}
Present the original lecture's definitions and notation before deriving results.

\section*{3. Interpretation}
Explain the assumptions, implications, and limitations at the required depth.

\end{document}
